\documentclass{beamer}
\usetheme{Madrid}
\usecolortheme{default}
\usepackage[utf8]{inputenc}
\usepackage[T1]{fontenc}
\usepackage[polish]{babel}
\usepackage{graphicx}
\usepackage{listings}
\usepackage{xcolor}

% Numery slajdów
\setbeamertemplate{footline}[frame number]

\title[Kompresja wag modeli SI]{Kompresja wag modeli sztucznej inteligencji \\ -- metody i wpływ}
\author{Jakub Grula}
\institute[KUL]{Katolicki Uniwersytet Lubelski Jana Pawła II \\ Wydział Filozofii \\ Instytut Filozofii}
\date{2026}

\begin{document}

\begin{frame}
    \titlepage
\end{frame}

\begin{frame}{Temat i cel pracy}
    \begin{block}{Problem}
        \begin{itemize}
            \item Modele SI są coraz większe (setki GB)
            \item Wysokie wymagania sprzętowe
            \item Ograniczenia wdrożeń produkcyjnych
            \item Brak popularnych rozwiązań dla TensorFlow
        \end{itemize}
    \end{block}
    
    \begin{block}{Cel pracy}
        \begin{itemize}
            \item Analiza metod kwantyzacji wag
            \item Implementacja własnego narzędzia \textit{Uni-Quant}
            \item Testy kompresji 4- i 8-bitowej
        \end{itemize}
    \end{block}
    
    \begin{alertblock}{Kwantyzacja}
        redukcja precyzji liczb (FP32 $\rightarrow$ INT4/INT8)
    \end{alertblock}
\end{frame}

\begin{frame}{Ograniczenia i wyzwania kwantyzacji}
    \begin{block}{Problemy}
        \begin{itemize}
            \item Utrata precyzji przy niskich bitach (4-bit)
            \item Wrażliwość na rozmiar bloku
            \item Wartości odstające w aktywacjach
        \end{itemize}
    \end{block}
    
    \begin{block}{Wyzwania}
        \begin{itemize}
            \item Dobór optymalnych parametrów
            \item Kompromis jakość/rozmiar
            \item Wsparcie dla różnych architektur
        \end{itemize}
    \end{block}
\end{frame}

\begin{frame}{Implementacja Uni-Quant}
    \begin{block}{Główne założenia}
        \begin{itemize}
            \item Blokowa kwantyzacja wag (rozmiar bloku: 8-128)
            \item Skalowanie per-blok: $\mathit{scale} = \frac{\max(|\mathit{block}|)}{2^{\mathit{bits}-1} - 1}$
            \item Pakowanie 2$\times$4-bit w 1 bajt
            \item Format archiwum \texttt{.uniq}
        \end{itemize}
    \end{block}
    
    \begin{alertblock}{Optymalizacja CUDA}
        kernele GPU przyspieszające kwantyzację i dekwantyzację
    \end{alertblock}
\end{frame}

\begin{frame}{Implementacja Uni-Quant c.d.}
    
    \begin{block}{Proces}
        \begin{enumerate}
            \item Iteracja po warstwach modelu TensorFlow
            \item Podział na bloki i obliczenie skal
            \item Kwantyzacja i zapis binarny
            \item Manifest JSON z metadanymi
        \end{enumerate}
    \end{block}
    
    \begin{exampleblock}{Uwagi}
        Małe fragmenty (< rozmiar bloku) zachowane w FP32
    \end{exampleblock}
\end{frame}

\begin{frame}{Wyniki testów}
    \begin{block}{Rodzaje testów}
        \begin{itemize}
            \item Regresja -- Niska strata na większych blokach
            \item Klasyfikacja binarna -- Niewielka strata na każdej konfiguracji
            \item Lasy losowe -- Brak straty (100\% przed i po)
            \item Klasyfikacja audio -- Mała strata na 8-bit z blokami 8-64
            \item Klasyfikacja obrazów -- Strata bliska zeru na 8-bit
        \end{itemize}
    \end{block}
    
    \begin{block}{Zależności}
        \begin{itemize}
            \item Mniej bitów $\rightarrow$ mniejszy rozmiar pliku
            \item Mniejszy blok $\rightarrow$ lepsza jakość, więcej metadanych
            \item Zachowanie FP32 dla małych warstw $\rightarrow$ mniejsza degradacja
        \end{itemize}
    \end{block}
\end{frame}

\end{document}
